\chapter{Kết luận và hướng phát triển}
\label{chap:conclusion}\label{ch:ketluan}

\section{Kết luận}

\begin{enumerate}
\item Xây dựng đường ống dữ liệu cho bốn bộ dữ liệu công khai (387 cell, 305.802 chu kỳ sau lọc) với bộ lọc ngoại lai nhân quả, và kiểm tra bằng số liệu vì sao dùng 16 đặc trưng sạc cùng chỉ số chu kỳ, vì sao bỏ các cách biến đổi khác (Mục~\ref{ch:dactrung}).
\item Thiết kế giao thức đánh giá không rò rỉ và đo tác động của các dạng rò rỉ: chia theo chu kỳ làm sai số thấp đi 1{,}2--1{,}8 lần, chuẩn hoá theo cả vòng đời cell làm sai số đổi 0{,}46--2{,}76 lần tuỳ bộ.
\item Đề xuất PINN bán giám sát với động học xám luôn không âm: ở 10\,\% nhãn sai số bằng 0{,}39--0{,}90 lần MLP, tiết kiệm 1{,}7--4{,}5 lần số nhãn; và mô hình kết hợp PINN-lai: 0{,}67--0{,}95 lần PINN-semi, thấp hơn GBDT ở 10/16 trường hợp.
\item So sánh có kiểm định với chín mô hình đối chứng, gồm bốn mô hình học sâu của đề cương (MLP, CNN, LSTM, Transformer), bằng MAE, RMSE, MAPE và $R^2$.
\item Chạy lại mã gốc PINN4SOH trên đúng 39 nhóm của bài (390 lần chạy bản công bố, 117 lần bản đã sửa), phát hiện loss đơn điệu của bản đã sinh số công bố không có gradient, và so cùng điều kiện: PINN-lai thấp hơn mã gốc chạy lại ở 12/12 chế độ, thấp hơn số công bố ở 11/12.
\item Đặt lại bài toán ít nhãn theo số lần đo dung lượng: đo thưa mọi cell với 1{,}2--6{,}0\,\% số chu kỳ cho sai số bằng 0{,}85--0{,}93 lần gắn nhãn đầy đủ cho 30\,\% số cell; mở rộng sang dự báo SOH, dung lượng còn lại và RUL.
\item Kết hợp GBDT có đặc trưng lịch sử SOH (mã của repo G-doan) với PINN-dự báo, đánh giá trong một giao thức đã sửa các điểm còn rò rỉ của E22 (E33, \Cref{sec:e33res}): khi có lịch sử đo dung lượng, GBDT lịch sử có MAE bằng 0{,}85--0{,}94 lần GBDT theo tầm, mô hình kết hợp mới bằng 0{,}79--0{,}99 lần kết hợp E22 và 1{,}01--1{,}48 lần GBDT lịch sử; với lịch sử ước lượng, kết hợp mới bằng 0{,}90--1{,}01 lần kết hợp E22. Theo quy tắc chọn sau thí nghiệm, với SOH, không mô hình nào thoả ở cả ba thiết lập, vì GBDT lịch sử có MAE trung bình thấp hơn từng mô hình còn lại trên ít nhất ba bộ nhưng kém có ý nghĩa trên TJU so với kết hợp mới và kết hợp E22; khi có lịch sử đo, GBDT lịch sử có MAE trung bình thấp nhất trên cả bốn bộ ở cả hai ngân sách; với RUL, quy tắc chọn kết hợp E22 ở cả hai kịch bản lịch sử.
\item Tổ chức lại mã nguồn thành các khối dữ liệu, mô hình, thí nghiệm, phân tích, so sánh với PINN4SOH, bản thảo bài báo và báo cáo; 11.429 lượt huấn luyện có kết quả trong báo cáo, mọi bảng sinh tự động từ kết quả, chạy lại cho kết quả trùng từng bit (Mục~\ref{sec:tailap}).
\end{enumerate}


\section{Hạn chế}

Hạn chế lớn nhất là một mô hình cây đơn giản (GBDT) vẫn là đối thủ khó nhất. PINN-semi đơn lẻ chưa vượt được nó khi có
từ 30\,\% nhãn, và chỉ khi ghép với GBDT thành PINN-lai mới thấp hơn GBDT ở phần lớn trường hợp. Ràng buộc vật lý chưa
được đưa vào chính mô hình cây. Phần vật lý cũng đóng góp chủ yếu qua tiên nghiệm đơn điệu và miền giá trị: các tham số
hoá lý $\lambda, E_a$ không được dữ liệu xác định, và tốc độ $r$ của dạng đạo hàm riêng không phải tốc độ suy giảm thật
của cell (\Cref{sec:vatly}). Đồ án không tái lập được chính xác số công bố của PINN4SOH vì không dựng lại được môi trường
gốc (GPU, phiên bản thư viện), nên báo cáo nêu cả số công bố lẫn số chạy lại. Khoảng dự đoán 90\,\% chưa đạt độ phủ danh
nghĩa ở mọi bộ.

Một số lựa chọn được đưa ra sau khi đã có kết quả trên tập kiểm tra: bản tinh chỉnh của PINN-lai, việc thêm đặc trưng
lịch sử trong HiPINN-lai, nhánh theo đời và dạng trung vị của HiPINN, và quy tắc chọn mô hình dự báo. Các kết quả này được
ghi rõ là phân tích hậu kiểm và chưa được kiểm chứng trên dữ liệu mới. Các lựa chọn rời rạc ở giai đoạn phát triển (thành
phần của PINN-lai, dạng đặc trưng lịch sử, độ dài cửa sổ của mô hình chuỗi) được chốt trên tập xác thực của lần lặp 0 gộp
qua 5 fold. Mỗi cell vì vậy có thể đã góp vào các lựa chọn đó khi nó là cell xác thực ở một fold khác, dù các lựa chọn
được chốt trước khi đọc dự đoán trên tập kiểm tra. Mọi tham số liên tục và mọi lựa chọn ở bước đánh giá đều làm riêng trong
từng fold. PINN-lai và các thí nghiệm về sau dùng 3 thay vì 5 lần lặp. Cuối cùng, đồ án chưa xử lý dữ liệu thô của từng bộ
nên chưa thử được các đặc trưng cần đường cong đầy đủ, và chưa có dữ liệu vận hành ngoài phòng thí nghiệm hay dữ liệu nhiều
cell ghép song song (giai đoạn 2 của đề cương).

\section{Hướng phát triển}

Hướng tiếp theo trước hết là giai đoạn 2 của đề cương: dữ liệu vận hành ngoài phòng thí nghiệm và dữ liệu nhiều cell
ghép song song. Trích lại đặc trưng từ tín hiệu điện áp, dòng điện gốc của từng bộ sẽ cho phép thử các đặc trưng cần
đường cong đầy đủ như dung lượng gia tăng hay điện áp hồi phục. Vì GBDT là đối thủ mạnh nhất của mạng có vật lý, một
hướng đáng thử là đưa ràng buộc vật lý vào chính mô hình cây. Khoảng dự đoán và khoảng RUL cần đạt độ phủ đặt ra, chẳng
hạn bằng cách hiệu chỉnh trên nhiều cell hơn hoặc theo từng cell. Các phân tích hậu kiểm và các lựa chọn mô hình cũng cần
được kiểm chứng trên một bộ dữ liệu chưa dùng khi phát triển.
